% 382_Myon_g2_Stand2025_De_ch.tex
% Johann Pascher, 29. Sept. 2026
% Myon g-2: Stand 2025 und Einordnung der FFGFT-Aussagen (Nachfolgedokument)

\providecommand{\KM}{\textbf{[K]}}
\providecommand{\SM}{\textbf{[S]}}

\chapter*{Myon \texorpdfstring{$g-2$}{g-2}: Stand 2025}
\addcontentsline{toc}{chapter}{Myon g-2: Stand 2025}

\begin{abstract}
Mehrere frühe Dokumente des Korpus beziehen sich auf die Abweichung zwischen dem
gemessenen anomalen magnetischen Moment des Myons und der Standardmodell-Vorhersage,
wie sie 2021 vorlag: $\Delta a_\mu = 251(59)\times10^{-11}$, rund $4{,}2\,\sigma$.
Mit dem Endergebnis des Fermilab-Experiments (Juni 2025) und dem White Paper~2025 der
Muon~$g-2$ Theory Initiative ist diese Abweichung auf $38(63)\times10^{-11}$
geschrumpft, also rund $0{,}6\,\sigma$ \KM. Aussagen, die einen T0-Zusatzbeitrag von
$251\times10^{-11}$ als Erklärung der Anomalie anführen, sind damit überholt; ein
solcher fester Beitrag läge heute rund $3{,}4\,\sigma$ über der Differenz \KM.

Die tragende $g-2$-Linie des Korpus (Dok.~018 Rev.~12, Dok.~158, A138) ist davon nicht
betroffen. Sie beruht auf \emph{Verhältnissen}: $\Delta a(\tau{-}\mu)/\Delta a(\mu{-}e)
=f^{1/3}-1$ und die Brücke $\Delta a(\tau{-}e)/\Delta a(\mu{-}e)=(144/125)\,m_\tau/m_\mu$.
In ihnen kürzen sich Projektionsfaktor und SI-Umrechnung heraus. Absolutwerte gewinnt
sie, indem sie die Verhältnisse an gemessenen Größen verankert; so folgt mit dem Stand
2025 $a_\tau=1{,}2811\times10^{-3}$ \KM. Dieser verankerte Wert ist der bessere
Absolutwert als der aus der direkten Formel, und er hängt nur an den Messwerten von
$a_e$ und $a_\mu$, nicht an der Standardmodell-Rechnung, deren Stand sich geändert hat.
\end{abstract}

\noindent\textbf{Zur Bezeichnung:} $f$ ist hier die Grundwindungszahl $f = 1/\xi = 7500$, keine Frequenz. In natürlichen Einheiten ist $f$ nur eine andere Lesart der Dualität $T\cdot m = 1$: $f = t_P/t_0 = \Lambda_{\text{T0}}/E_P$ mit $t_0 = \xi\,t_P$ und $\Lambda_{\text{T0}} = E_P/\xi$, und aus $\xi E_0^2 = 1$ folgt $f = E_0^2$. Als dimensionsloses Verhältnis lässt sich $f$ nicht begründet auf~1 setzen (Dok.~A135).

% ============================================================
\section*{Statusmarkierungen}
% ============================================================
\begin{center}
\begin{tabular}{@{}lp{10cm}@{}}
\textbf{[K]} & Numerisch bestätigt (Prüfskript). \\[3pt]
\textbf{[S]} & Spekulativ: Hinweis oder Vermutung, im Korpus noch nicht hergeleitet oder geprüft. \\
\end{tabular}
\end{center}
Die experimentellen und theoretischen Eingangswerte stammen aus den im
Literaturverzeichnis genannten Quellen; sie werden hier nicht hergeleitet, sondern
übernommen.

% ============================================================
\section{Ausgangslage im Korpus}
% ============================================================
Im Korpus gibt es zwei Arten von Aussagen zum $g-2$ der Leptonen, die man
auseinanderhalten muss.

\textbf{(a) Erklärung der Anomalie durch einen T0-Zusatzbeitrag.} Frühe Dokumente
setzen einen Beitrag $\Delta a_\mu^{(\mathrm{T0})}=251\times10^{-11}$ an oder
bezeichnen die Myon-Anomalie als durch den $\xi$-Term erklärt bzw. als
Vergleichsmaßstab: Dok.~019, 030, 033, 049, 053, 059, 070, 081, 095 sowie in der
A-Serie A138 (Vergleich mit der „Myon-Diskrepanz``). Diese Aussagen hängen direkt an
der Größe der Abweichung zwischen Messung und Standardmodell.

\textbf{(b) Die Verhältnis-Linie.} Dok.~018 (Rev.~12), Dok.~158 und A138 beschreiben
die Leptonen durch fraktale Windungsexponenten ($5/3$ für das Myon, $4/3$ für das Tau)
und stellen die Verhältnisse in den Mittelpunkt: Verhältnisse von $g-2$-Differenzen
sind strukturell exakt, weil sich Projektionsfaktor, fraktale Korrektur und
SI-Umrechnung herauskürzen. Absolutwerte folgen erst durch Verankerung an einer
gemessenen Größe -- dieselbe Arbeitsteilung wie bei den Massen, wo Verhältnisse
geometrisch festliegen und ein empirischer Anker die Skala setzt (A138, Dok.~351).
Diese Linie vergleicht mit Messwerten, nicht mit der Differenz zum Standardmodell.

% ============================================================
\section{Stand 2025}
% ============================================================
\begin{itemize}
\item \textbf{Messung.} Das Fermilab-Experiment E989 hat im Juni 2025 sein
  Endergebnis veröffentlicht: $a_\mu=116\,592\,070{,}5(14{,}8)\times10^{-11}$,
  eine Genauigkeit von $127$~ppb \cite{FNAL2025}. Zusammen mit E821 (Brookhaven)
  ergibt sich das Weltmittel $a_\mu^{\mathrm{exp}}=116\,592\,071{,}5(14{,}5)\times10^{-11}$
  \cite{WP25}.
\item \textbf{Theorie.} Das White Paper~2025 der Muon $g-2$ Theory Initiative gibt
  $a_\mu^{\mathrm{SM}}=116\,592\,033(62)\times10^{-11}$ an \cite{WP25}. Der
  wesentliche Unterschied zu 2020 ist die hadronische Vakuumpolarisation: Sie wird
  jetzt aus Gitter-QCD bestimmt statt aus $e^+e^-$-Wirkungsquerschnitten.
\item \textbf{Differenz.} $a_\mu^{\mathrm{exp}}-a_\mu^{\mathrm{SM}}=38(63)\times10^{-11}$,
  also rund $0{,}6\,\sigma$ \KM. 2021 waren es $251(59)\times10^{-11}$, rund
  $4{,}2\,\sigma$ \KM.
\item \textbf{Offen innerhalb der Theorie.} Warum die datengetriebene Bestimmung der
  hadronischen Vakuumpolarisation von der Gitter-Bestimmung abweicht, ist noch nicht
  geklärt \cite{CERNCourier2025}. Die Anomalie gilt damit als stark geschwächt, aber
  nicht endgültig abgeschlossen.
\end{itemize}
Geändert hat sich also die Standardmodell-Rechnung, nicht der Messwert: Das Weltmittel
von $a_\mu$ hat sich seit 2021 nur um rund $10\times10^{-11}$ verschoben.

% ============================================================
\section{Folgen für die Aussagen vom Typ (a)}
% ============================================================
Ein fester T0-Zusatzbeitrag von $251\times10^{-11}$ war 2021 genau so groß wie die
Abweichung, die er erklären sollte. Heute liegt er um
\[
  \frac{251-38}{63}\approx3{,}4\,\sigma
\]
über der gemessenen Differenz \KM. Die Aussage „der $\xi$-Term erklärt die
Myon-Anomalie`` ist damit überholt. Für einen zusätzlichen Beitrag jenseits des
Standardmodells bleibt nach heutigem Stand bei zwei Standardabweichungen ein
Spielraum von etwa $-88$ bis $+164\times10^{-11}$ \KM. Ein T0-Beitrag in dieser
Größenordnung ist nicht ausgeschlossen, er ist aber im Korpus nicht hergeleitet
\SM.

Betroffen sind die in Abschnitt~1 unter (a) genannten Dokumente. Sie verwenden die
Abweichung von 2021 nicht mehr als Beleg oder Zielgröße und verweisen für die
Einordnung der Myon-Anomalie auf dieses Dokument.

% ============================================================
\section{Die Verhältnis-Linie ist nicht betroffen}
% ============================================================
Der Kern von Dok.~018, Dok.~158 und A138 sind zwei Verhältnisaussagen \KM:
\[
  \frac{\Delta a(\tau{-}\mu)}{\Delta a(\mu{-}e)} = f^{1/3}-1 = 18{,}57,
  \qquad
  \frac{\Delta a(\tau{-}e)}{\Delta a(\mu{-}e)} = \frac{144}{125}\,\frac{m_\tau}{m_\mu} = 19{,}37 .
\]
Die erste folgt allein aus den Exponenten $5/3$ und $4/3$. Die zweite ersetzt
$f^{1/3}=19{,}57$ durch das Verhältnis der Massenformeln und verbindet so $g-2$ mit
den Leptonmassen; $f$, $k_\mathrm{geom}$ und $K_\mathrm{frak}$ kürzen sich heraus.

\textbf{Vom Verhältnis zum Absolutwert.} Verankert man die Brücke an den gemessenen
Werten von $a_e$ und $a_\mu$, ergibt sich der Absolutwert für das Tau:
\[
  a_\tau = a_e + \frac{144}{125}\,\frac{m_\tau}{m_\mu}\,(a_\mu-a_e) = 1{,}2811\times10^{-3}
\]
mit dem Stand 2025 \KM; Dok.~018 und A138 nennen $1{,}282\times10^{-3}$ mit gerundetem
$a_e$, Dok.~158 $1{,}281\times10^{-3}$. Die direkte Absolutformel
$a_\tau=a_e+4\pi/f^{4/3}$ ergibt dagegen $1{,}2454\times10^{-3}$ \KM. Der verankerte
Wert ist der bessere, weil er nur das exakte Verhältnis und gemessene Größen enthält;
deshalb verwenden Dok.~158 und A138 ihn als Vorhersage.

\textbf{Warum die Anomalie hier keine Rolle spielt.} Der Anker ist die gemessene
Differenz $a_\mu-a_e=6{,}2685\times10^{-6}$. Sie hat sich zwischen 2021 und 2025 nur um
rund $10^{-10}$ geändert \KM. Die Anomalie war eine Aussage über die
Standardmodell-Rechnung; die Verhältnis-Linie benutzt diese Rechnung nicht. Die
Tau-Vorhersage bleibt deshalb unverändert. Sie liegt um $1{,}04\times10^{-4}$ über der
Standardmodell-Vorhersage $a_\tau^{\mathrm{SM}}=1{,}1772\times10^{-3}$ \KM. Die
heutigen experimentellen Schranken für $a_\tau$ sind rund 85-mal weiter als dieser
Unterschied (CMS 2024: $-0{,}0042 < a_\tau < 0{,}0046$, 95\,\% CL) \KM; die Entscheidung steht aus, bis $a_\tau$ gemessen ist \SM.

% ============================================================
\section{Die direkten Absolutformeln}
% ============================================================
Neben der Verhältnis-Linie enthält Dok.~018 direkte Absolutformeln aus der
Windungsgeometrie, korrigiert um $K_\mathrm{frak}^{3/2}$. Mit den Werten aus Dok.~018
ergibt sich \KM:
\begin{align*}
  a_e^{(\mathrm{korr})} &= 1{,}1598\times10^{-3}\quad(0{,}014\,\%\ \text{vom Messwert}),\\
  a_\mu^{(\mathrm{korr})} &= 1{,}1642\times10^{-3}\quad(-0{,}15\,\%\ \text{vom Weltmittel}).
\end{align*}
Auch diese Werte vergleichen mit Messwerten und hängen nicht an der Anomalie. Sie
arbeiten aber auf dem Niveau von $10^{-4}$ bis $10^{-3}$ und sind kein
Präzisionsvergleich dort, wo die Anomalie diskutiert wurde. Ein Zahlenbeispiel zeigt,
warum der Korpus für Absolutwerte den Weg über Verhältnisse wählt: Die direkt
berechnete Differenz $\Delta a(\mu{-}e)=4\pi/f^{5/3}=4{,}373\times10^{-6}$ liegt rund
$30\,\%$ unter der gemessenen \KM, während das Verhältnis der Differenzen exakt aus den
Exponenten folgt. Der verankerte Weg umgeht diese Unsicherheit.

\textbf{A138.} Der dortige Vergleich, die Tau-Abweichung sei „siebenmal größer als die
Myon-Diskrepanz``, ist gegenstandslos, weil es keine signifikante Myon-Diskrepanz mehr
gibt. Die Tau-Aussage selbst ist davon unabhängig.

% ============================================================
\section{Ergebnis}
% ============================================================
Die Myon-$g-2$-Anomalie von 2021 ist mit dem Stand 2025 nicht mehr signifikant.
FFGFT-Aussagen, die sie als Beleg oder als Zielgröße eines T0-Zusatzbeitrags
verwenden, sind überholt; die betroffenen Dokumente verweisen dafür auf dieses Dokument. Die tragende
Verhältnis-Linie ist nicht betroffen: Ihre Verhältnisse sind strukturell exakt, ihre
Absolutwerte gewinnt sie durch Verankerung an gemessenen Größen, und die Messwerte
haben sich nicht geändert. Ihre Tau-Vorhersage $a_\tau\approx1{,}281\times10^{-3}$
bleibt die eigentliche Testaussage. Ob die FFGFT darüber hinaus einen Beitrag im heute
erlaubten Bereich von etwa $10^{-9}$ zum Myon liefert, ist offen \SM.

\medskip\noindent Prüfskript: \texttt{2/python/Dok382\_Skripte/pruef\_382\_myon\_g2\_stand.py}, 21/21.

\begin{thebibliography}{9}
\bibitem{FNAL2025} Muon $g-2$ Collaboration (D.~P.~Aguillard et al.),
  \emph{Measurement of the Positive Muon Anomalous Magnetic Moment to 127~ppb},
  Phys.\ Rev.\ Lett.\ 135 (2025) 101802, arXiv:2506.03069.
\bibitem{WP25} R.~Aliberti et al. (Muon $g-2$ Theory Initiative),
  \emph{The anomalous magnetic moment of the muon in the Standard Model: an update},
  Phys.\ Rep.\ 1143 (2025) 1--158, arXiv:2505.21476.
\bibitem{CERNCourier2025} CERN Courier, \emph{Fermilab's final word on muon g-2} (2025).
\end{thebibliography}
